\documentclass{article}
\usepackage{pathfinder-readable}

\title{When a Contract Leaves Cooperation Unfinished}

\begin{document}
\maketitle

\begin{abstract}
This note connects a theory of incentives for AI agents with a study of language model agents making contracts
in a two-player game. The thread asked whether the weaker results for natural-language contracts reveal a
problem with the theory's rule about verifiable capabilities. That link did not hold, but a closer question
emerged: whether an incomplete contract discourages the further cooperation needed to finish. The peers
established a path-based lower bound on that remaining need, while the reported game outcomes suggest, but do
not prove, such a discouragement effect. The verifier paused the thread because the accepted obligations and
game histories needed to test the explanation were not available.
\end{abstract}

\section{Introduction}

Bergemann, Koh and Morris study how a designer can induce an AI agent to report what it can do and then follow
instructions \cite{Q}. An agent may hide evidence of a capability, but cannot manufacture evidence it lacks.
The paper separates actions an agent can physically take from actions it has an incentive to take. It also
asks how much a reward can depend on what the designer can observe. Its main mechanism results concern a
static setting, even though the authors discuss the difficulty of extending it to long sequences of actions.

Wyse, Bustos, Volkova and Kleiman-Weiner study language model agents on a coloured board \cite{P}. Each player
wants to reach a goal but starts without chips of the other player's colour. Players can arrange to pay for
each other's moves and can negotiate a binding tile contract before play. In one format, the agreed tile
obligations appear as structured data and are checked by code. In another, the negotiation remains in prose
and a judge decides whether each proposed move is covered. A promised transfer occurs only if its giver has
the chip.

The thread began with the thought that prose interpretation might violate the first paper's rule against
counterfeit capability claims. The peers rejected that link: the rule governs evidence for reports about an
agent's type, whereas the judge in the board game interprets an agreement about a later move. The more useful
connection is the distinction between a feasible plan and incentives to carry it out, together with the
question of what an agreement's enforcement can recognise.

\section{What was done}

The peers first checked what the board-game results actually say. On boards where each player needs the
other's help, both players finish in \(46\%\) of games with natural-language tile contracts, \(79\%\) with
programmatic tile contracts, and \(61\%\) without contracts \cite{P}. The two contract formats cover about
\(2.6\) tiles per game on average. Yet players proposed later pay-for-partner arrangements less often with
prose contracts than with programmatic ones, \(1.4\) rather than \(3.0\) proposals per game, and accepted
fewer, \(0.7\) rather than \(1.4\). When declining to trade, players cited an existing contract more often in
the prose setting, \(94\%\) rather than \(74\%\) of those refusals. The peers took these observations as a
reason to examine whether an incomplete contract crowds out later cooperation.

They then replaced the average number of covered tiles with a test for each board and contract. Let \(i\)
denote one player, \(C\) the accepted tile obligations, and \(\mathcal P_i\) the simple paths from that
player's start to its goal. For a tile \(v\), an obligation in \(C\) \emph{covers} \(i\) at \(v\) if it pays
for that player's move onto \(v\). They defined the uncovered path requirement
\[
\delta_i(C)=\min_{p\in\mathcal P_i}
\#\{v\in p:\ v\text{ has the other player's colour and }C
\text{ does not cover }i\text{ at }v\}.
\]
Here \(p\) is one path and \(\#\) counts its qualifying tile visits. This is a lower bound on the further help
player \(i\) needs, not a forecast of which path the agents will choose. The peers next checked why the bound
cannot by itself certify that a contract suffices. Finally, they proposed comparing the same obligations shown
as structured data or as faithfully generated prose, with identical enforcement, while checking enforcement
separately.

\section{Results}
\enlargethispage{3\baselineskip}

The path bound holds under the game's stated chip and contract rules \cite{P}. A player begins with no chips
of the other colour. Any successful walk contains a simple start-to-goal path. On every such path, each
other-colour tile not covered by the contract must be paid for through later voluntary help. Thus a player who
finishes needs at least \(\delta_i(C)\) such payments or transfers. In particular, \(\delta_i(C)>0\) proves
that the contract alone cannot get that player to the goal. The two peers independently checked this argument
in the ledger, with the qualification that it counts a necessary minimum, not a sufficient plan.

The reverse implication fails. Even when \(\delta_i(C)=0\), the giver may have spent a promised chip when the
covered move occurs. The two players' individually suitable routes may also be incompatible in timing or
inventory. A stronger physical test would search joint states containing both positions, both inventories,
remaining obligations and turn order, asking whether both goals can be reached with no further voluntary
exchange. Such a search would establish possibility only. It would not show that either player prefers to
follow that route. This is the practical form of the feasibility-versus-obedience distinction in the first
paper \cite{Q}.

\section{Objections and unresolved questions}

The initial claim about counterfeit capability reports failed because no such reports are shown in the
board-game evidence. The judge observes the proposed destination; the uncertain issue is whether the prose
agreement covers it. The first paper's discussion of imperfectly measured, payoff-relevant conditions offers a
way to think about that issue, but gives no predicted ranking between the two contract formats \cite{Q}.

The board-game paper audits \(1{,}155\) approved prose-contract moves and finds three judge errors, all
mistaken approvals \cite{P}. This weakens an account based on frequent mistaken approvals. It does not measure
mistaken denials, and a genuinely ambiguous clause may have no single correct reading until its meaning is
fixed independently. Similar average numbers of covered tiles also do not show that the formats promise the
same tiles or cover equally useful paths.

The lower bound has not been calculated for the observed contracts. The supplied paper and thread contain
neither the accepted obligations and game histories needed for that calculation nor an audit of denied judge
decisions. The peers therefore did not establish that unfinished obligations caused the completion gap. Even a
gap under matched obligations and enforcement could reflect changed beliefs, planning or a choice among
equilibria. The first paper's static results, with broad reward instruments, do not settle play in this
sequential game with limited contract terms \cite{Q,P}.

\section{Why the thread stopped}

The verifier accepted the path bound but paused the thread because the proposed explanation remained a
research design, not a supported finding about this pair. The smallest restart is to obtain the accepted
contracts and game histories, translate each contract into tile obligations, calculate \(\delta_i(C)\), and
compare later trading and completion by the resulting shortfall. A joint reachability check would sharpen the
physical classification. If the claim is specifically that prose presentation causes the gap, the next
experiment must hold the board, obligations, information and enforcement fixed while varying only their
display; a separate audit must include denied as well as approved judge decisions.

\bibliographystyle{plainurl}
\bibliography{references}

\end{document}
